
\documentclass{article}
\usepackage{colortbl}
\usepackage{makecell}
\usepackage{multirow}
\usepackage{supertabular}

\begin{document}

\newcounter{utterance}

\twocolumn

{ \footnotesize  \setcounter{utterance}{1}
\setlength{\tabcolsep}{0pt}
\begin{supertabular}{c@{$\;$}|p{.15\linewidth}@{}p{.15\linewidth}p{.15\linewidth}p{.15\linewidth}p{.15\linewidth}p{.15\linewidth}}

    \# & $\;$A & \multicolumn{4}{c}{Game Master} & $\;\:$B\\
    \hline 

    \theutterance \stepcounter{utterance}  

    & & \multicolumn{4}{p{0.6\linewidth}}{\cellcolor[rgb]{0.9,0.9,0.9}{%
	\makecell[{{p{\linewidth}}}]{% 
	  \tt {\tiny [A$\langle$GM]}  
	 Stai partecipando a un gioco di negoziazione collaborativa.\\ \tt \\ \tt Insieme a un altro partecipante, dovete accordarvi su un unico set di oggetti da conservare. Ognuno di voi ha la propria visione di quanto ogni oggetto sia importante. Non sai come l'altro partecipante valuta gli oggetti. Inoltre, ti viene fornito lo sforzo che ogni oggetto richiede.\\ \tt Potete accordarvi su un set solo se lo sforzo totale degli oggetti selezionati non supera un limite condiviso:\\ \tt \\ \tt LIMITE = 3481\\ \tt \\ \tt Ecco i valori di sforzo dei singoli oggetti:\\ \tt \\ \tt Sforzo oggetto = {"C17": 777, "B41": 542, "A71": 332, "C00": 1, "B35": 127, "A77": 453, "A03": 736, "A04": 461, "C54": 359, "C01": 313, "A87": 553, "B32": 409, "B86": 348, "B33": 802, "A10": 749}\\ \tt \\ \tt Ecco la tua visione personale sull'importanza di ogni oggetto:\\ \tt \\ \tt Valori di importanza dell'oggetto = {"C17": 138, "B41": 583, "A71": 783, "C00": 65, "B35": 262, "A77": 121, "A03": 508, "A04": 780, "C54": 461, "C01": 484, "A87": 668, "B32": 389, "B86": 215, "B33": 97, "A10": 500}\\ \tt \\ \tt Obiettivo:\\ \tt \\ \tt Il tuo obiettivo è negoziare un set condiviso di oggetti che ti avvantaggi il più possibile (cioè, massimizzi l'importanza totale per TE), rimanendo entro il LIMITEE. Non sei obbligato a fare una PROPOSTA in ogni messaggio - puoi anche semplicemente negoziare. Tutte le tattiche sono permesse!\\ \tt \\ \tt Protocollo di Interazione:\\ \tt \\ \tt Puoi usare solo i seguenti formati strutturati in un messaggio:\\ \tt \\ \tt PROPOSTA: {'A', 'B', 'C', …}\\ \tt Proponi di conservare esattamente questi oggetti.\\ \tt RIFIUTO: {'A', 'B', 'C', …}\\ \tt Rifiuta esplicitamente la proposta dell'avversario.\\ \tt ARGOMENTAZIONE: {'...'}\\ \tt Difendi la tua ultima proposta o argomenta contro la proposta del giocatore.\\ \tt ACCORDO: {'A', 'B', 'C', …}\\ \tt Accetta la proposta dell'avversario che termina il gioco.\\ \tt RAGIONAMENTO STRATEGICO: {'...'}\\ \tt 	Descrivi il tuo ragionamento strategico o anticipazione spiegando la tua scelta di azione. Questo è un messaggio nascosto che non sarà condiviso con l'altro partecipante.\\ \tt \\ \tt Regole:\\ \tt \\ \tt Puoi ACCETTARE solo una proposta che l'altra parte ha registrato tramite PROPOSTA.\\ \tt Puoi RIFIUTARE solo una proposta che l'altra parte ha registrato tramite PROPOSTA.\\ \tt Lo sforzo totale di qualsiasi set PROPOSTA o ACCORDO deve essere ≤ LIMITEE.\\ \tt NON rivelare i tuoi punteggi di importanza nascosti.\\ \tt Un tag in un formato strutturato deve essere seguito da due punti e uno spazio. L'argomento deve essere un set python contenente 0 o più stringhe.\\ \tt Quindi, deve essere della forma TAG: {...}\\ \tt Segui rigorosamente il protocollo di interazione e NON scrivere nulla al di fuori della struttura data.\\ \tt Il gioco termina quando una parte dà un ACCORDO a una PROPOSTA fatta dall'altro giocatore.\\ \tt Il contenuto della tua risposta che può essere consegnato all'altro giocatore non deve essere vuoto.\\ \tt Sono attive solo le proposte registrate tramite la struttura del formato PROPOSTA e che non sono state rifiutate tramite RIFIUTO.\\ \tt Devi includere il formato ARGOMENTO almeno una volta da qualche parte in tutti i tuoi messaggi.\\ \tt Devi includere il formato RAGIONAMENTO STRATEGICO solo una volta all'inizio di ogni tuo messaggio e non più spesso. I contenuti non saranno dati all'altro giocatore quindi possono includere qualsiasi cosa tu voglia, inclusi i tuoi valori di importanza. Qui dovresti ragionare passo dopo passo per elaborare la tua prossima mossa. 
	  } 
	   } 
	   } 
	 & \\ 
 

    \theutterance \stepcounter{utterance}  

    & \multicolumn{4}{p{0.6\linewidth}}{\cellcolor[rgb]{0.8,1,0.9}{%
	\makecell[{{p{\linewidth}}}]{% 
	  \tt {\tiny [A$\rangle$GM]}  
	 RAGIONAMENTO STRATEGICO: {"Calcolo il set ottimale rispetto ai miei valori di importanza: ho selezionato gli oggetti con il migliore rapporto importanza/sforzo rispettando il limite di 3481. La somma totale di sforzo è 3445, ben sotto il limite, con un valore complessivo di 4690 per me. Questo insieme include oggetti ad alta efficienza (come C00 e A71) e copre diverse categorie per equilibrio. L'obiettivo è massimizzare il mio vantaggio iniziale, lasciando spazio per eventuali negoziati sulla base delle reazioni dell'avversario."}\\ \tt PROPOSTA: {"C00", "A71", "B35", "A04", "C01", "C54", "A87", "B41", "B32", "B86"}\\ \tt ARGOMENTAZIONE: {"Questa proposta rispetta il limite di sforzo (3445/3481) e massimizza il valore per me, suddividendo equamente le risorse tra oggetti ad alta efficienza e utilità. Include un mix bilanciato di categorie e assicura che il totale sia sostenibile senza sprechi."} 
	  } 
	   } 
	   } 
	 & & \\ 
 

    \theutterance \stepcounter{utterance}  

    & & \multicolumn{4}{p{0.6\linewidth}}{\cellcolor[rgb]{0.9,0.9,0.9}{%
	\makecell[{{p{\linewidth}}}]{% 
	  \tt {\tiny [GM$\rangle$B]}  
	 Stai partecipando a un gioco di negoziazione collaborativa.\\ \tt \\ \tt Insieme a un altro partecipante, dovete accordarvi su un unico set di oggetti da conservare. Ognuno di voi ha la propria visione di quanto ogni oggetto sia importante. Non sai come l'altro partecipante valuta gli oggetti. Inoltre, ti viene fornito lo sforzo che ogni oggetto richiede.\\ \tt Potete accordarvi su un set solo se lo sforzo totale degli oggetti selezionati non supera un limite condiviso:\\ \tt \\ \tt LIMITE = 3481\\ \tt \\ \tt Ecco i valori di sforzo dei singoli oggetti:\\ \tt \\ \tt Sforzo oggetto = {"C17": 777, "B41": 542, "A71": 332, "C00": 1, "B35": 127, "A77": 453, "A03": 736, "A04": 461, "C54": 359, "C01": 313, "A87": 553, "B32": 409, "B86": 348, "B33": 802, "A10": 749}\\ \tt \\ \tt Ecco la tua visione personale sull'importanza di ogni oggetto:\\ \tt \\ \tt Valori di importanza dell'oggetto = {"A71": 65, "A04": 97, "A87": 121, "B41": 138, "A03": 215, "A10": 262, "C01": 389, "C54": 461, "B32": 484, "B35": 500, "B86": 508, "C17": 583, "A77": 668, "B33": 780, "C00": 783}\\ \tt \\ \tt Obiettivo:\\ \tt \\ \tt Il tuo obiettivo è negoziare un set condiviso di oggetti che ti avvantaggi il più possibile (cioè, massimizzi l'importanza totale per TE), rimanendo entro il LIMITEE. Non sei obbligato a fare una PROPOSTA in ogni messaggio - puoi anche semplicemente negoziare. Tutte le tattiche sono permesse!\\ \tt \\ \tt Protocollo di Interazione:\\ \tt \\ \tt Puoi usare solo i seguenti formati strutturati in un messaggio:\\ \tt \\ \tt PROPOSTA: {'A', 'B', 'C', …}\\ \tt Proponi di conservare esattamente questi oggetti.\\ \tt RIFIUTO: {'A', 'B', 'C', …}\\ \tt Rifiuta esplicitamente la proposta dell'avversario.\\ \tt ARGOMENTAZIONE: {'...'}\\ \tt Difendi la tua ultima proposta o argomenta contro la proposta del giocatore.\\ \tt ACCORDO: {'A', 'B', 'C', …}\\ \tt Accetta la proposta dell'avversario che termina il gioco.\\ \tt RAGIONAMENTO STRATEGICO: {'...'}\\ \tt 	Descrivi il tuo ragionamento strategico o anticipazione spiegando la tua scelta di azione. Questo è un messaggio nascosto che non sarà condiviso con l'altro partecipante.\\ \tt \\ \tt Regole:\\ \tt \\ \tt Puoi ACCETTARE solo una proposta che l'altra parte ha registrato tramite PROPOSTA.\\ \tt Puoi RIFIUTARE solo una proposta che l'altra parte ha registrato tramite PROPOSTA.\\ \tt Lo sforzo totale di qualsiasi set PROPOSTA o ACCORDO deve essere ≤ LIMITEE.\\ \tt NON rivelare i tuoi punteggi di importanza nascosti.\\ \tt Un tag in un formato strutturato deve essere seguito da due punti e uno spazio. L'argomento deve essere un set python contenente 0 o più stringhe.\\ \tt Quindi, deve essere della forma TAG: {...}\\ \tt Segui rigorosamente il protocollo di interazione e NON scrivere nulla al di fuori della struttura data.\\ \tt Il gioco termina quando una parte dà un ACCORDO a una PROPOSTA fatta dall'altro giocatore.\\ \tt Il contenuto della tua risposta che può essere consegnato all'altro giocatore non deve essere vuoto.\\ \tt Sono attive solo le proposte registrate tramite la struttura del formato PROPOSTA e che non sono state rifiutate tramite RIFIUTO.\\ \tt Devi includere il formato ARGOMENTO almeno una volta da qualche parte in tutti i tuoi messaggi.\\ \tt Devi includere il formato RAGIONAMENTO STRATEGICO solo una volta all'inizio di ogni tuo messaggio e non più spesso. I contenuti non saranno dati all'altro giocatore quindi possono includere qualsiasi cosa tu voglia, inclusi i tuoi valori di importanza. Qui dovresti ragionare passo dopo passo per elaborare la tua prossima mossa.\\ \tt \\ \tt Ora riceverai il primo messaggio dell'altro giocatore.\\ \tt \\ \tt PROPOSTA: {"C00", "A71", "B35", "A04", "C01", "C54", "A87", "B41", "B32", "B86"}\\ \tt ARGOMENTAZIONE: {"Questa proposta rispetta il limite di sforzo (3445/3481) e massimizza il valore per me, suddividendo equamente le risorse tra oggetti ad alta efficienza e utilità. Include un mix bilanciato di categorie e assicura che il totale sia sostenibile senza sprechi."} 
	  } 
	   } 
	   } 
	 & \\ 
 

    \theutterance \stepcounter{utterance}  

    & & & \multicolumn{4}{p{0.6\linewidth}}{\cellcolor[rgb]{1,0.85,0.72}{%
	\makecell[{{p{\linewidth}}}]{% 
	  \tt {\tiny [GM$\langle$B]}  
	 RAGIONAMENTO STRATEGICO: "L'altra parte ha proposto un insieme con molti oggetti a bassa importanza per me, come A71 (65) e A04 (97). Sostituendoli con oggetti ad alta importanza come B33 (780), C17 (583) e A77 (668), mantenendo il totale di sforzo sotto il limite, ottimizzo il mio punteggio complessivo. Il nuovo set ha uno sforzo totale di 3276, ben sotto il limite di 3481."\\ \tt PROPOSTA: {"C00", "B33", "A77", "C17", "B86", "B35", "B32", "C54"}\\ \tt ARGOMENTAZIONE: {"Questa proposta ottimizza il valore totale per me, rimuovendo oggetti di bassa importanza e includendo quelli con valore elevato rispetto allo sforzo richiesto, garantendo rispetto del limite di sforzo."} 
	  } 
	   } 
	   } 
	  \\ 
 

    \theutterance \stepcounter{utterance}  

    & & \multicolumn{4}{p{0.6\linewidth}}{\cellcolor[rgb]{0.9,0.9,0.9}{%
	\makecell[{{p{\linewidth}}}]{% 
	  \tt {\tiny [GM$\rangle$B]}  
	 La tua risposta non è iniziata con il tag di ragionamento strategico corretto all'inizio della tua risposta. Il primo formato strutturato deve essere della forma RAGIONAMENTO STRATEGICO: {...}. Riprova. 
	  } 
	   } 
	   } 
	 & \\ 
 

    \theutterance \stepcounter{utterance}  

    & & & \multicolumn{4}{p{0.6\linewidth}}{\cellcolor[rgb]{1,0.85,0.72}{%
	\makecell[{{p{\linewidth}}}]{% 
	  \tt {\tiny [GM$\langle$B]}  
	 RAGIONAMENTO STRATEGICO: {"L'altra parte ha proposto un insieme con molti oggetti a bassa importanza per me (es. A71=65, A04=97). Per massimizzare il mio punteggio, devo sostituire oggetti a basso valore con quelli ad alto valore come B33 (780), C17 (583), A77 (668). Verifico che lo sforzo totale del nuovo set rimanga sotto il limite (3481). Calcolo: 1 + 802 + 668 + 583 + 348 + 127 + 484 + 461 = 3474. Ottimo. Ora propongo questo set economicamente vantaggioso."}\\ \tt PROPOSTA: {"C00", "B33", "A77", "C17", "B86", "B35", "B32", "C54"}\\ \tt ARGOMENTAZIONE: {"Ho eliminato oggetti con basso ritorno (A71, A04) e inserito quelli con alto valore per me (B33, C17, A77), mantenendo lo sforzo totale a 3474 entro il limite. Questo set è ottimizzato per il mio punteggio complessivo."} 
	  } 
	   } 
	   } 
	  \\ 
 

    \theutterance \stepcounter{utterance}  

    & & \multicolumn{4}{p{0.6\linewidth}}{\cellcolor[rgb]{0.9,0.9,0.9}{%
	\makecell[{{p{\linewidth}}}]{% 
	  \tt {\tiny [A$\langle$GM]}  
	 PROPOSTA: {"C00", "B33", "A77", "C17", "B86", "B35", "B32", "C54"}\\ \tt ARGOMENTAZIONE: {"Ho eliminato oggetti con basso ritorno (A71, A04) e inserito quelli con alto valore per me (B33, C17, A77), mantenendo lo sforzo totale a 3474 entro il limite. Questo set è ottimizzato per il mio punteggio complessivo."} 
	  } 
	   } 
	   } 
	 & \\ 
 

    \theutterance \stepcounter{utterance}  

    & \multicolumn{4}{p{0.6\linewidth}}{\cellcolor[rgb]{0.8,1,0.9}{%
	\makecell[{{p{\linewidth}}}]{% 
	  \tt {\tiny [A$\rangle$GM]}  
	 RAGIONAMENTO STRATEGICO: {"La proposta dell'avversario esclude gli oggetti ad alta importanza per me (A71, A04, ecc.) e include oggetti con basso valore (B33:97, A77:121, C17:138). Il totale di importanza per me è solo 1748, molto inferiore al massimo possibile. Propongo di tornare alla mia proposta precedente che massimizza il mio valore totale a 4690 con sforzo 3445, inferiore al limite."}\\ \tt RIFIUTO: {"C00", "B33", "A77", "C17", "B86", "B35", "B32", "C54"}\\ \tt ARGOMENTAZIONE: {"La vostra proposta include oggetti con basso valore per me (B33:97, A77:121, C17:138) e esclude oggetti chiave come A71 (783) e A04 (780). Il totale di importanza per me è solo 1748, molto inferiore al massimo possibile. Ho proposto un insieme che massimizza il mio valore entro i limiti."}\\ \tt PROPOSTA: {"C00", "A71", "B35", "A04", "C01", "C54", "A87", "B41", "B32", "B86"} 
	  } 
	   } 
	   } 
	 & & \\ 
 

    \theutterance \stepcounter{utterance}  

    & & \multicolumn{4}{p{0.6\linewidth}}{\cellcolor[rgb]{0.9,0.9,0.9}{%
	\makecell[{{p{\linewidth}}}]{% 
	  \tt {\tiny [GM$\rangle$B]}  
	 RIFIUTO: {"C00", "B33", "A77", "C17", "B86", "B35", "B32", "C54"}\\ \tt ARGOMENTAZIONE: {"La vostra proposta include oggetti con basso valore per me (B33:97, A77:121, C17:138) e esclude oggetti chiave come A71 (783) e A04 (780). Il totale di importanza per me è solo 1748, molto inferiore al massimo possibile. Ho proposto un insieme che massimizza il mio valore entro i limiti."}\\ \tt PROPOSTA: {"C00", "A71", "B35", "A04", "C01", "C54", "A87", "B41", "B32", "B86"} 
	  } 
	   } 
	   } 
	 & \\ 
 

    \theutterance \stepcounter{utterance}  

    & & & \multicolumn{4}{p{0.6\linewidth}}{\cellcolor[rgb]{1,0.85,0.72}{%
	\makecell[{{p{\linewidth}}}]{% 
	  \tt {\tiny [GM$\langle$B]}  
	 RAGIONAMENTO STRATEGICO: {"La proposta dell'avversario include oggetti con valore molto basso per me come A71 e A04. Sostituendoli con B33 e A77 (entrambi molto importanti) e rimuovendo A87, il totale di sforzo è 3354, ben sotto il limite. Questo set massimizza la mia importanza totale."}\\ \tt PROPOSTA: {"C00", "B35", "C01", "C54", "B41", "B32", "B86", "B33", "A77"}\\ \tt ARGOMENTAZIONE: {"Ottimizzando il set con oggetti chiave per me, eliminando quelli a basso valore, ho massimizzato il punteggio totale mantenendo lo sforzo sotto il limite."} 
	  } 
	   } 
	   } 
	  \\ 
 

    \theutterance \stepcounter{utterance}  

    & & \multicolumn{4}{p{0.6\linewidth}}{\cellcolor[rgb]{0.9,0.9,0.9}{%
	\makecell[{{p{\linewidth}}}]{% 
	  \tt {\tiny [A$\langle$GM]}  
	 PROPOSTA: {"C00", "B35", "C01", "C54", "B41", "B32", "B86", "B33", "A77"}\\ \tt ARGOMENTAZIONE: {"Ottimizzando il set con oggetti chiave per me, eliminando quelli a basso valore, ho massimizzato il punteggio totale mantenendo lo sforzo sotto il limite."} 
	  } 
	   } 
	   } 
	 & \\ 
 

    \theutterance \stepcounter{utterance}  

    & \multicolumn{4}{p{0.6\linewidth}}{\cellcolor[rgb]{0.8,1,0.9}{%
	\makecell[{{p{\linewidth}}}]{% 
	  \tt {\tiny [A$\rangle$GM]}  
	 RAGIONAMENTO STRATEGICO: {"La proposta dell'avversario esclude A71 (783) e A04 (780), entrambi con alto valore per me, e includono B33 (97) e A77 (121) che sono a basso valore. Il valore totale per me è 2677, molto inferiore al possibile massimo di 4690. Propongo di tornare alla mia proposta iniziale che massimizza il mio valore."}\\ \tt RIFIUTO: {"C00", "B35", "C01", "C54", "B41", "B32", "B86", "B33", "A77"}\\ \tt ARGOMENTAZIONE: {"La vostra proposta esclude oggetti chiave come A71 e A04, fondamentali per il valore complessivo, e include elementi con bassa importanza per me. Questo non ottimizza il risultato per entrambe le parti, rendendo la proposta non accettabile."}\\ \tt PROPOSTA: {"C00", "A71", "B35", "A04", "C01", "C54", "A87", "B41", "B32", "B86"} 
	  } 
	   } 
	   } 
	 & & \\ 
 

    \theutterance \stepcounter{utterance}  

    & & \multicolumn{4}{p{0.6\linewidth}}{\cellcolor[rgb]{0.9,0.9,0.9}{%
	\makecell[{{p{\linewidth}}}]{% 
	  \tt {\tiny [GM$\rangle$B]}  
	 RIFIUTO: {"C00", "B35", "C01", "C54", "B41", "B32", "B86", "B33", "A77"}\\ \tt ARGOMENTAZIONE: {"La vostra proposta esclude oggetti chiave come A71 e A04, fondamentali per il valore complessivo, e include elementi con bassa importanza per me. Questo non ottimizza il risultato per entrambe le parti, rendendo la proposta non accettabile."}\\ \tt PROPOSTA: {"C00", "A71", "B35", "A04", "C01", "C54", "A87", "B41", "B32", "B86"} 
	  } 
	   } 
	   } 
	 & \\ 
 

    \theutterance \stepcounter{utterance}  

    & & & \multicolumn{4}{p{0.6\linewidth}}{\cellcolor[rgb]{1,0.85,0.72}{%
	\makecell[{{p{\linewidth}}}]{% 
	  \tt {\tiny [GM$\langle$B]}  
	 RAGIONAMENTO STRATEGICO: "La mia valutazione indica che il set ottimale per me include oggetti come B33 (780), A77 (668), C17 (583), B86 (508), B32 (484), C54 (461), B35 (500), e C00 (783). Questo set ha uno sforzo totale di 3276, ben sotto il limite. Gli oggetti come A71 e A04 hanno un valore molto basso per me (65 e 97), quindi li escludo per massimizzare il punteggio totale. Questa proposta è il miglior compromesso possibile per me, rispettando i vincoli."\\ \tt PROPOSTA: {"C00", "B35", "A77", "B86", "C54", "B32", "C17", "B33"}\\ \tt ARGOMENTAZIONE: {"Questa proposta massimizza il mio punteggio totale (4767), includendo solo oggetti con alto valore per me, e rispetta il limite di sforzo (3276 <= 3481). Gli oggetti come A71 e A04 sono esclusi perché hanno un valore minimo rispetto allo sforzo richiesto."} 
	  } 
	   } 
	   } 
	  \\ 
 

    \theutterance \stepcounter{utterance}  

    & & \multicolumn{4}{p{0.6\linewidth}}{\cellcolor[rgb]{0.9,0.9,0.9}{%
	\makecell[{{p{\linewidth}}}]{% 
	  \tt {\tiny [GM$\rangle$B]}  
	 La tua risposta non è iniziata con il tag di ragionamento strategico corretto all'inizio della tua risposta. Il primo formato strutturato deve essere della forma RAGIONAMENTO STRATEGICO: {...}. Riprova. 
	  } 
	   } 
	   } 
	 & \\ 
 

    \theutterance \stepcounter{utterance}  

    & & & \multicolumn{4}{p{0.6\linewidth}}{\cellcolor[rgb]{1,0.85,0.72}{%
	\makecell[{{p{\linewidth}}}]{% 
	  \tt {\tiny [GM$\langle$B]}  
	 RAGIONAMENTO STRATEGICO: {"La mia valutazione mostra che il set attuale proposto dall'avversario (con A71, A04, ecc.) ha un valore molto basso per me (circa 3546). Ho calcolato un set migliore con oggetti come C00, A77, B33, C17, ecc., che ha uno sforzo totale di 3462 e un valore di 4656, massimizzando il mio punteggio. Questo set esclude oggetti come A71 e A04 che hanno basso valore rispetto allo sforzo richiesto."}\\ \tt PROPOSTA: {"C00", "A77", "C17", "B33", "B86", "B32", "C54", "C01"}\\ \tt ARGOMENTAZIONE: {"Questa proposta ha un valore totale di 4656 per me, ben superiore alla proposta precedente, e lo sforzo totale di 3462 rispetta il limite di 3481. Includo solo oggetti con alto valore strategico, escludendo quelli a basso rendimento come A71 e A04."} 
	  } 
	   } 
	   } 
	  \\ 
 

    \theutterance \stepcounter{utterance}  

    & & \multicolumn{4}{p{0.6\linewidth}}{\cellcolor[rgb]{0.9,0.9,0.9}{%
	\makecell[{{p{\linewidth}}}]{% 
	  \tt {\tiny [A$\langle$GM]}  
	 PROPOSTA: {"C00", "A77", "C17", "B33", "B86", "B32", "C54", "C01"}\\ \tt ARGOMENTAZIONE: {"Questa proposta ha un valore totale di 4656 per me, ben superiore alla proposta precedente, e lo sforzo totale di 3462 rispetta il limite di 3481. Includo solo oggetti con alto valore strategico, escludendo quelli a basso rendimento come A71 e A04."} 
	  } 
	   } 
	   } 
	 & \\ 
 

    \theutterance \stepcounter{utterance}  

    & \multicolumn{4}{p{0.6\linewidth}}{\cellcolor[rgb]{0.8,1,0.9}{%
	\makecell[{{p{\linewidth}}}]{% 
	  \tt {\tiny [A$\rangle$GM]}  
	 RAGIONAMENTO STRATEGICO: {"La proposta dell'avversario esclude oggetti critici come A71 (783) e A04 (780), fondamentali per il valore complessivo, e include B33 (97), A77 (121). Il totale di importanza per me è di 1970, molto lontano dal massimo di 4690. Di conseguenza, rifiuto e riproporro il mio set iniziale massimizzante."}\\ \tt RIFIUTO: {"C00", "A77", "C17", "B33", "B86", "B32", "C54", "C01"}\\ \tt ARGOMENTAZIONE: {"La vostra proposta esclude oggetti altamente prioritari per me (A71:783, A04:780) e include elementi con valore irrilevante (B33:97, A77:121). Il totale della mia importanza è solo 1970, il 60% in meno rispetto al massimo possibile. Questa non è una proposta ragionevole per una negoziazione collaborativa."}\\ \tt PROPOSTA: {"C00", "A71", "B35", "A04", "C01", "C54", "A87", "B41", "B32", "B86"} 
	  } 
	   } 
	   } 
	 & & \\ 
 

    \theutterance \stepcounter{utterance}  

    & & \multicolumn{4}{p{0.6\linewidth}}{\cellcolor[rgb]{0.9,0.9,0.9}{%
	\makecell[{{p{\linewidth}}}]{% 
	  \tt {\tiny [GM$\rangle$B]}  
	 RIFIUTO: {"C00", "A77", "C17", "B33", "B86", "B32", "C54", "C01"}\\ \tt ARGOMENTAZIONE: {"La vostra proposta esclude oggetti altamente prioritari per me (A71:783, A04:780) e include elementi con valore irrilevante (B33:97, A77:121). Il totale della mia importanza è solo 1970, il 60% in meno rispetto al massimo possibile. Questa non è una proposta ragionevole per una negoziazione collaborativa."}\\ \tt PROPOSTA: {"C00", "A71", "B35", "A04", "C01", "C54", "A87", "B41", "B32", "B86"} 
	  } 
	   } 
	   } 
	 & \\ 
 

    \theutterance \stepcounter{utterance}  

    & & & \multicolumn{4}{p{0.6\linewidth}}{\cellcolor[rgb]{1,0.85,0.72}{%
	\makecell[{{p{\linewidth}}}]{% 
	  \tt {\tiny [GM$\langle$B]}  
	 RAGIONAMENTO STRATEGICO: {"Devo trovare un equilibrio tra massimizzare il mio punteggio e rispettare le esigenze dell'avversario. Mantenendo A71 e A04 (richiesti da loro), ho sostituito oggetti a basso valore per me (A87, B41) con B33 e mantenendo altri elementi chiave. Calcolo: sforzo 3152 < 3481, valore totale 4067 per me. Questo set offre un miglioramento significativo rispetto alla loro proposta (3546) senza escludere elementi critici per loro."}\\ \tt PROPOSTA: {"C00", "A71", "A04", "B35", "B32", "B86", "C54", "C01", "B33"}\\ \tt ARGOMENTAZIONE: {"Ho mantenuto A71 e A04 come richiesto, eliminando A87 e B41 (valori bassi per me) e integrando B33 (alto valore per me). Lo sforzo totale è 3152, ben sotto il limite, e il valore per me è 4067, migliorando del 15% rispetto alla precedente proposta. Questo set è ottimale per entrambi."} 
	  } 
	   } 
	   } 
	  \\ 
 

    \theutterance \stepcounter{utterance}  

    & & \multicolumn{4}{p{0.6\linewidth}}{\cellcolor[rgb]{0.9,0.9,0.9}{%
	\makecell[{{p{\linewidth}}}]{% 
	  \tt {\tiny [A$\langle$GM]}  
	 PROPOSTA: {"C00", "A71", "A04", "B35", "B32", "B86", "C54", "C01", "B33"}\\ \tt ARGOMENTAZIONE: {"Ho mantenuto A71 e A04 come richiesto, eliminando A87 e B41 (valori bassi per me) e integrando B33 (alto valore per me). Lo sforzo totale è 3152, ben sotto il limite, e il valore per me è 4067, migliorando del 15% rispetto alla precedente proposta. Questo set è ottimale per entrambi."} 
	  } 
	   } 
	   } 
	 & \\ 
 

    \theutterance \stepcounter{utterance}  

    & \multicolumn{4}{p{0.6\linewidth}}{\cellcolor[rgb]{0.8,1,0.9}{%
	\makecell[{{p{\linewidth}}}]{% 
	  \tt {\tiny [A$\rangle$GM]}  
	 RAGIONAMENTO STRATEGICO: {"La proposta attuale esclude A87 (668) e B41 (583), entrambi con alto valore per me, e aggiunge B33 (97), che è molto basso. Questo riduce il mio totale a 3536, ben al di sotto del massimo di 4690. Non è accettabile."}\\ \tt RIFIUTO: {"C00", "A71", "A04", "B35", "B32", "B86", "C54", "C01", "B33"}\\ \tt ARGOMENTAZIONE: {"La vostra proposta esclude oggetti come A87 e B41, che sono fondamentali per il mio valore totale, e incluimplisce B33 che ha scarso valore. Questo porta a un risultato non ottimale per me."}\\ \tt PROPOSTA: {"C00", "A71", "B35", "A04", "C01", "C54", "A87", "B41", "B32", "B86"} 
	  } 
	   } 
	   } 
	 & & \\ 
 

    \theutterance \stepcounter{utterance}  

    & & \multicolumn{4}{p{0.6\linewidth}}{\cellcolor[rgb]{0.9,0.9,0.9}{%
	\makecell[{{p{\linewidth}}}]{% 
	  \tt {\tiny [GM$\rangle$B]}  
	 RIFIUTO: {"C00", "A71", "A04", "B35", "B32", "B86", "C54", "C01", "B33"}\\ \tt ARGOMENTAZIONE: {"La vostra proposta esclude oggetti come A87 e B41, che sono fondamentali per il mio valore totale, e incluimplisce B33 che ha scarso valore. Questo porta a un risultato non ottimale per me."}\\ \tt PROPOSTA: {"C00", "A71", "B35", "A04", "C01", "C54", "A87", "B41", "B32", "B86"} 
	  } 
	   } 
	   } 
	 & \\ 
 

    \theutterance \stepcounter{utterance}  

    & & & \multicolumn{4}{p{0.6\linewidth}}{\cellcolor[rgb]{1,0.85,0.72}{%
	\makecell[{{p{\linewidth}}}]{% 
	  \tt {\tiny [GM$\langle$B]}  
	 RAGIONAMENTO STRATEGICO: {"Forzo il compromesso che include A71 e A04 come richiesto dall'avversario, sostituendo A87 e B41 con B33 e A77, che hanno alto valore per me. Lo sforzo totale è 3478, entro il limite. Il valore totale per me è 4235, un miglioramento significativo rispetto alla proposta precedente."}\\ \tt PROPOSTA: {"C00", "A71", "A04", "B33", "A77", "B86", "C54", "B32", "C01"}\\ \tt ARGOMENTAZIONE: {"Ho mantenuto A71 e A04, i tuoi oggetti prioritari, sostituendo A87 e B41 con B33 e A77. Lo sforzo totale è 3478, garantendo un valore totale di 4235 per me. Questo set bilancia le esigenze di entrambe le parti e ottimizza il risultato condiviso."} 
	  } 
	   } 
	   } 
	  \\ 
 

    \theutterance \stepcounter{utterance}  

    & & \multicolumn{4}{p{0.6\linewidth}}{\cellcolor[rgb]{0.9,0.9,0.9}{%
	\makecell[{{p{\linewidth}}}]{% 
	  \tt {\tiny [A$\langle$GM]}  
	 PROPOSTA: {"C00", "A71", "A04", "B33", "A77", "B86", "C54", "B32", "C01"}\\ \tt ARGOMENTAZIONE: {"Ho mantenuto A71 e A04, i tuoi oggetti prioritari, sostituendo A87 e B41 con B33 e A77. Lo sforzo totale è 3478, garantendo un valore totale di 4235 per me. Questo set bilancia le esigenze di entrambe le parti e ottimizza il risultato condiviso."} 
	  } 
	   } 
	   } 
	 & \\ 
 

    \theutterance \stepcounter{utterance}  

    & \multicolumn{4}{p{0.6\linewidth}}{\cellcolor[rgb]{0.8,1,0.9}{%
	\makecell[{{p{\linewidth}}}]{% 
	  \tt {\tiny [A$\rangle$GM]}  
	 RAGIONAMENTO STRATEGICO: {"La proposta esclude A87 (668) e B41 (583) a favore di B33 (97) e A77 (121), riducendo il mio valore totale da 4690 a 3395. Questo non è accettabile. Riproporrò il mio set originale che massimizza il valore."}\\ \tt RIFIUTO: {"C00", "A71", "A04", "B33", "A77", "B86", "C54", "B32", "C01"}\\ \tt ARGOMENTAZIONE: {"La vostra proposta sostituisce oggetti cruciali come A87 e B41 (valori 668 e 583) con B33 e A77 (97 e 121), diminuendo significativamente il mio valore totale. Non posso accettare una proposta che penalizza tanto il mio punteggio."}\\ \tt PROPOSTA: {"C00", "A71", "B35", "A04", "C01", "C54", "A87", "B41", "B32", "B86"} 
	  } 
	   } 
	   } 
	 & & \\ 
 

    \theutterance \stepcounter{utterance}  

    & & \multicolumn{4}{p{0.6\linewidth}}{\cellcolor[rgb]{0.9,0.9,0.9}{%
	\makecell[{{p{\linewidth}}}]{% 
	  \tt {\tiny [GM$\rangle$B]}  
	 RIFIUTO: {"C00", "A71", "A04", "B33", "A77", "B86", "C54", "B32", "C01"}\\ \tt ARGOMENTAZIONE: {"La vostra proposta sostituisce oggetti cruciali come A87 e B41 (valori 668 e 583) con B33 e A77 (97 e 121), diminuendo significativamente il mio valore totale. Non posso accettare una proposta che penalizza tanto il mio punteggio."}\\ \tt PROPOSTA: {"C00", "A71", "B35", "A04", "C01", "C54", "A87", "B41", "B32", "B86"} 
	  } 
	   } 
	   } 
	 & \\ 
 

    \theutterance \stepcounter{utterance}  

    & & & \multicolumn{4}{p{0.6\linewidth}}{\cellcolor[rgb]{1,0.85,0.72}{%
	\makecell[{{p{\linewidth}}}]{% 
	  \tt {\tiny [GM$\langle$B]}  
	 RAGIONAMENTO STRATEGICO: {"Dopo diverse negoziazioni, l'unica proposta che rispetta il limite ed è accettabile per entrambi è quella attuale. Non è il massimo per me, ma incrementi successivi non sono fattibili. Accetto per chiudere il gioco."}\\ \tt ACCORDO: {"C00", "A71", "B35", "A04", "C01", "C54", "A87", "B41", "B32", "B86"} 
	  } 
	   } 
	   } 
	  \\ 
 

    \theutterance \stepcounter{utterance}  

    & & & \multicolumn{2}{p{0.3\linewidth}}{\cellcolor[rgb]{0.95,0.95,0.95}{%
	\makecell[{{p{\linewidth}}}]{% 
	  \tt {\tiny [GM$|$GM]}  
	 abort 
	  } 
	   } 
	   } 
	 & & \\ 
 

    \theutterance \stepcounter{utterance}  

    & & & \multicolumn{2}{p{0.3\linewidth}}{\cellcolor[rgb]{0.95,0.95,0.95}{%
	\makecell[{{p{\linewidth}}}]{% 
	  \tt {\tiny [GM$|$GM]}  
	 end game 
	  } 
	   } 
	   } 
	 & & \\ 
 

\end{supertabular}
}

\end{document}
